\section{Gaussian moments at sampled edges}\label{sec:entropy}

We use the spectral estimates of Proposition~\ref{prop:spectral-input}
to identify the empirical distribution of finitely many eigenvector
coordinates at an edge. The estimate must be quantitative: a fixed
moment will later be summed over a subpolynomial number of eigenindices.
We develop the entropy framework of \citet{BackhauszSzegedy}
quantitatively to obtain the conditional moment accuracy required here.

\begin{proposition}[Uniform endpoint moments]\label{prop:edge-moments}
Fix positive integers \(k,\ell\) and a nonnegative integer \(q\).
There are constants \(c,C>0\) and events of graphs of probability tending
to one with the following property. Choose distinct indices
\(i(1),\ldots,i(\ell)\) satisfying \(|x_{i(j)}|\le S^2\), by any rule
depending only on the graph. Independently sample \(k\) uniform oriented
edges \((v_\alpha,v'_\alpha)\), conditionally on the graph, and set
\[
 b_j=\bigl(u_{i(j)}(v_1),\ldots,u_{i(j)}(v_k),
           u_{i(j)}(v'_1),\ldots,u_{i(j)}(v'_k)\bigr)^\top .
\]
For every coordinate monomial \(P\) of total degree at most \(q\),
\[
 \left|\EE\bigl[P(b_1,\ldots,b_\ell)\mid A\bigr]
       -\EE P(Z_1,\ldots,Z_\ell)\right|\le Cn^{-c},
 \qquad Z_1,\ldots,Z_\ell\ \text{independent }N(0,C_k).
\]
The constants and graph events are uniform over the indicated choices
of eigenindices. They may depend on \(d,E,k,\ell,q\); no assertion is
made for dimensions or moment orders growing with \(n\).
\end{proposition}

We first prove the assertion for one sampled edge. Write
\[
 V(v)=\bigl(u_{i(1)}(v),\ldots,u_{i(\ell)}(v)\bigr)\in\R^\ell.
\]
A star observation consists of the value at a uniform vertex followed
by its \(d\) neighbor values in uniform random order. An edge observation
uses a uniform oriented edge. Throughout this section the graph is
conditioned on, and randomness in an observation includes its location.
Write \(H\) for discrete entropy and \(\mathcal H\) for differential
entropy, both with natural logarithms. Positive polynomial error
exponents may decrease from one occurrence to the next.

The proof passes from colored-label counts to an entropy inequality
for Gaussian-smoothed stars and edges. We remove the nearly
deterministic eigenvector relation, then use tensor coercivity to turn
the resulting entropy bound into Gaussian moment estimates.

\subsection{Counting colored stars}

We begin with a counting statement independent of eigenvectors.
For a labeling of the vertices by an alphabet of size \(M\), let \(T\)
be the probability distribution of the ordered star observation,
symmetrized over all permutations of the neighbor positions. Let \(Q\)
be the oriented edge distribution.
The counting argument follows the colored-star configuration-model
framework of \citet[Section~4]{BackhauszSzegedyTrees}; see also the
asymptotic expected-coloring formula of
\citet[Theorem~2.3]{BackhauszBordenaveSzegedy}.
We retain explicit errors for an alphabet growing with \(n\).

\begin{lemma}[Uniform type count]\label{ent:type-count}
Suppose \(M=M_n\) and, for some fixed \(\kappa>0\),
\(M^{d+1}\log n\le n^{1-\kappa}\).
With probability tending to one, simultaneously for every pair of
histograms \(T,Q\), the number of vertex labelings with these histograms
is at most
\begin{equation}\label{eq:ent-type-count}
 \exp\left\{nH(T)-\frac{nd}{2}H(Q)+O(n^{1-c})\right\},
\end{equation}
where \(c>0\) and the error constant depend only on \(d,\kappa\).
\end{lemma}

\begin{proof}
Work first in the configuration model, with an ordering of the \(d\)
half-edges at each vertex. Prescribe the vertex labels and the ordered
list of desired neighbor labels at every vertex. For a fixed ordered
tuple type \(T_0\), there are at most \(\exp(nH(T_0))\) such arrays, by
the multinomial bound.

Let \(m_{jk}=ndQ_{jk}\) be the number of half-edges whose vertex has
color \(j\) and whose prescribed neighbor has color \(k\).
Compatibility requires \(m_{jk}=m_{kj}\), and \(m_{jj}\) must be even.
The probability that a uniform matching respects all these prescriptions is
\[
 \frac{\displaystyle\prod_{j<k}m_{jk}!
                    \prod_j(m_{jj}-1)!!}{(nd-1)!!},
\]
with zero-size factors omitted. Stirling's bounds give its logarithm as
\[
 -\frac{nd}{2}H(Q)+O(M^2\log n).
\]
Thus the expected number of labelings of ordered type \(T_0\) has the
claimed exponential bound with \(H(T_0)\) in place of \(H(T)\).

Symmetrization increases entropy, so \(H(T_0)\le H(T)\); it leaves \(Q\)
unchanged. There are at most
\(\exp(O(M^{d+1}\log n))\) possible ordered tuple types. Summing over
those with a given symmetrization therefore changes only the sublinear
error. Markov's inequality, followed by a union bound over all types,
gives \eqref{eq:ent-type-count} with a slightly larger sublinear
allowance. The simplicity probability is bounded away from zero by
Lemma~\ref{lem:model}, so the statement transfers to uniform simple
graphs. The symmetrized histograms no longer depend on the half-edge
ordering.
\end{proof}

\subsection{A quantitative differential-entropy inequality}

Choose small fixed parameters \(\gamma,\alpha>0\) satisfying
\begin{equation}\label{eq:ent-parameters}
 2\ell(d+1)\gamma<1,\qquad
 4\ell\gamma+\alpha<\frac14,\qquad
 0<\alpha<\gamma,
\end{equation}
and take \(\alpha\) smaller than the rigidity exponent in
Proposition~\ref{prop:spectral-input}. Put \(\eta=n^{-\alpha}\).
Given the graph, let \(Y\) have \(\ell\) independent Gaussian colors,
each with covariance
\begin{equation}\label{eq:ent-noise}
 \Gamma=\frac{1}{\pi\rho}\,
          \frac{\eta}{(A-E)^2+\eta^2}.
\end{equation}
For \(0\le\epsilon\le1\), set \(F=\epsilon V+Y\).
The parameter \(\eta\) regularizes the eigenvector constraint; the
independent parameter \(\epsilon\) controls the size of the eigenvector
signal. We retain the entropy inequality for every \(\epsilon\), then
choose a small power of \(n^{-1}\) when extracting each fixed moment.
All eigenvectors used here satisfy the uniform subpolynomial
coordinate bound from Proposition~\ref{prop:spectral-input}.

\begin{lemma}[Star--edge entropy]\label{ent:star-edge}
On events of graphs of probability tending to one, uniformly over the
choices of \(V\) and \(0\le\epsilon\le1\),
\begin{equation}\label{eq:ent-star-edge}
 \mathcal H(F_{\rm star})-\frac d2\mathcal H(F_{\rm edge})
 \ge \frac1n\mathcal H(Y_{\rm global})-O(n^{-c}),
\end{equation}
where \(Y_{\rm global}\) is the complete \(n\ell\)-dimensional array.
Moreover,
\begin{equation}\label{eq:ent-global-entropy}
 \frac1n\mathcal H(Y_{\rm global})
 =\ell\left(\frac12\log\frac{2\pi{\rm e}\eta}{\pi\rho}
                    -U(E)\right)+O(n^{-c}),
 \qquad
 U(E)=\int\log|t-E|\rho_d(t)\,\dd t .
\end{equation}
\end{lemma}

\begin{proof}
Quantize each coordinate into intervals of width \(b=n^{-\gamma}\),
clipping all values outside a range of order
\([-n^\gamma,n^\gamma]\) into its two end bins. The resulting vertex
alphabet has \(M=n^{2\ell\gamma+o(1)}\) elements, so
Lemma~\ref{ent:type-count} applies. Let \(\mathcal L\) be the random
quantized labeling of \(F\), and let \(T,Q\) be its random histograms.
Entropy of the type index, followed by entropy bounded by logarithmic
type size, gives
\begin{equation}\label{eq:ent-quantized}
 \frac1nH(\mathcal L)
 \le H(\EE T)-\frac d2\EE H(Q)+O(n^{-c}).
\end{equation}
Here concavity was used only for the positive star term.

We next replace \(\EE H(Q)\) by \(H(\EE Q)\). Two independent edge
locations have Gaussian observations whose off-diagonal covariance
block has average Frobenius norm at most
\[
 C\left(\frac{\Tr\Gamma^2}{n^2}\right)^{1/2}
 \le \frac{C}{\sqrt n\,\eta}.
\]
Indeed endpoints from different sampled edges are independent uniform
vertices, and \(\Tr\Gamma^2\le Cn/\eta^2\).
Each individual edge covariance has smallest eigenvalue at least
\(c_1\eta\), since \(\|A\|\le d\).
After whitening the two marginal blocks, Gaussian total variation
from the product law is bounded by the off-block norm when that norm
is small, and by one otherwise. The small-norm bound follows, for
example, by differentiating the Gaussian density along the straight
line joining the two covariance matrices. Averaging yields
\[
 \EE_{\rm locations}\,
 \dist_{\TV}\bigl(\Law(Y_{\rm edge},Y_{\rm edge}'),
                 \Law(Y_{\rm edge})\otimes
                 \Law(Y_{\rm edge}')\bigr)
 \le \frac{C}{\sqrt n\,\eta^2},
\]
where the marginal laws on the right are at the fixed sampled
locations. Deterministic shifts and subsequent quantization do not
increase this distance.

Consequently each cell frequency of \(Q\) has variance
\(O(n^{-1/2}\eta^{-2})\). There are \(M^2\) edge cells, so
\[
 \EE\|Q-\EE Q\|_1
 \le n^{4\ell\gamma-1/4+\alpha+o(1)}=O(n^{-c}).
\]
The finite-alphabet entropy continuity bound
\(H_{\rm bin}(h)+h\log(M^2)\) at total variation distance \(h\)
now proves
\begin{equation}\label{eq:ent-edge-concentration}
 \EE H(Q)=H(\EE Q)+O(n^{-c}).
\end{equation}

We record the differential-entropy passage carefully. Every fixed
location covariance, including the complete array covariance, has
spectrum in \([c_1\eta,C/\eta]\); every coordinate of the deterministic
shift is \(n^{o(1)}\). Unbounded binning and the entropy bound on a cube give
\[
 \frac1nH(\mathcal L)
 \ge \frac1n\mathcal H(Y_{\rm global})
           -\ell\log b-O(n^{-c}).
\]
The displayed error accounts for clipping. It is in fact
superpolynomially small: one-coordinate Gaussian tails beyond
\(n^\gamma\), together with their quantized entropy contributions,
are superpolynomially small uniformly in the shifts, because
\(\alpha<\gamma\). The loss for the full array is at most the sum
of the coordinate losses. Translation by the deterministic array
\(\epsilon V\) does not change its differential entropy.

For either fixed-dimensional local observation \(X=F_{\rm star}\)
or \(F_{\rm edge}\), we also have
\begin{equation}\label{eq:ent-binning}
 H(\text{binned }X)
 =\mathcal H(X)-(\dim X)\log b+O(n^{-c}).
\end{equation}
To justify the two-sided estimate, fill each bin uniformly.
Bin averaging a Gaussian density changes it in total variation by
at most \(Cb\eta^{-1/2}\), by integrating its absolute first
derivatives; the same bound holds for a mixture. Clipping adds only
a superpolynomial tail error. Both resulting densities have
polynomial supremum and second-moment bounds.
Such bounds convert polynomial total variation into polynomial
entropy error: subtract the common minimum of the two densities.
The remaining densities have equal mass \(h\); after division by
\(h\), their supremum and second moment are bounded by a polynomial
divided by \(h\), so their entropies have absolute value
\(O(1+\log(n/h))\). The entropy errors in mixing them with the
common part are at most the binary entropy. This proves
\eqref{eq:ent-binning}, and will also apply to fixed-dimensional
projections below.

The laws \(\EE T,\EE Q\) are exactly the binned local laws.
Substitute \eqref{eq:ent-edge-concentration} and
\eqref{eq:ent-binning} into \eqref{eq:ent-quantized}.
The mesh terms cancel, since
\((d+1)\ell-(d/2)(2\ell)=\ell\).
This proves \eqref{eq:ent-star-edge}.

Finally the Gaussian determinant formula gives
\[
 \frac1n\mathcal H(Y_{\rm global})
 =\frac{\ell}{2}\log\frac{2\pi{\rm e}\eta}{\pi\rho}
       -\frac{\ell}{n}\sum_i\log|\lambda_i-E-i\eta|.
\]
The test function in this sum has Lipschitz constant \(O(\eta^{-1})\).
Rigidity therefore gives polynomial error in replacing the eigenvalues
by Kesten--McKay quantiles. The quantile sum differs from its integral
by \(O((n\eta)^{-1})\), by summing the lengths of the quantile
intervals. An omitted or extra bounded eigenvalue costs
\(O(\log n/n)\). Lemma~\ref{lem:tree-identities} gives the remaining
regularization error \(O(\eta)\). This proves
\eqref{eq:ent-global-entropy}.
\end{proof}

\subsection{Removing the approximate eigenvector constraint}

The entropy inequality contains one nearly deterministic direction:
the sum of the neighbor values is nearly \(E\) times the center.
We remove that direction so that the limiting Gaussian has identity
covariance. Write a star as \((Z,X)\), with
\(Z\in\R^\ell\) and \(X\in\R^d\otimes\R^\ell\), and let \(P\) be
orthogonal projection of \(\R^d\) onto \({\bf1}^{\perp}\).
Define
\[
 \mathcal B(Z,X)=(Z,PX/\sqrt a),\qquad
 R_{\rm res}(Z,X)=\sum_{i=1}^dX_i-EZ,\qquad
 D=(\R\oplus{\bf1}^{\perp})\otimes\R^\ell.
\]
Here and below projections on spatial coordinates are tensored with
the color identity. Let \(w\) denote the unit central direction in
\(\R\oplus{\bf1}^{\perp}\), and set
\[
 v_i=pw+\sqrt a\,Pe_i,\qquad
 U_i=\spanop(w,v_i)\otimes\R^\ell .
\]
Then \(\|v_i\|=1\) and \(\ip{w}{v_i}=p\). Let \(P_i\) be orthogonal
projection onto \(U_i\), put \(B_V=\mathcal B(V_{\rm star})\), and
let \(G\) be an independent standard Gaussian in \(D\). Entropies on
these spaces use their Euclidean volume.

\begin{lemma}[Projected-star entropy]\label{ent:projected-entropy}
On the graph events above, uniformly for \(0\le\epsilon\le1\) and
each \(i\),
\begin{equation}\label{eq:ent-projected-entropy}
 \mathcal H(\epsilon B_V+G)
       -\frac d2\mathcal H(P_i(\epsilon B_V+G))
 \ge -O(n^{-c}).
\end{equation}
Moreover \(B_V\) has mean zero and second-moment matrix
\(\Id_D+O(n^{-1+o(1)})\). Every fixed moment of \(P_iB_V\) is
invariant up to \(O(n^{-1+o(1)})\) under the orthogonal map
interchanging \(w\) and \(v_i\), acting trivially on the colors.
\end{lemma}

\begin{proof}
The absolute Jacobian determinant of
\((Z,X)\mapsto(\mathcal B(Z,X),R_{\rm res}(Z,X))\) is
\((\sqrt d\,a^{-(d-1)/2})^\ell\).
For each color, the residual variance of \(Y\), averaged over
centers, is at most \(\eta/(\pi\rho)\), since
\[
 (A-E)\Gamma(A-E)
 =\frac{\eta}{\pi\rho}
   \frac{(A-E)^2}{(A-E)^2+\eta^2}
 \le\frac{\eta}{\pi\rho}\Id.
\]
The eigenvector equations give a residual
\(O(n^{-1+o(1)})\) for \(V\), uniformly over locations and choices:
\(|\lambda_{i(j)}-E|\le S^2/n\).
More precisely, put
\(\Delta=\diag(\lambda_{i(j)}-E:1\le j\le\ell)\).
At a uniform center \(v\), the deterministic residual is
\(\epsilon\Delta V(v)\), whose second-moment matrix is
\(\epsilon^2\Delta^2\) by eigenvector orthogonality.
The Gaussian colors are independent and centered conditionally on
the location, so cross terms vanish and
\begin{equation}\label{eq:ent-residual-covariance}
 \Cov\bigl(R_{\rm res}(F_{\rm star})\bigr)
 \le \frac{\eta}{\pi\rho}\Id_\ell+\epsilon^2\Delta^2
 \le \left(\frac{\eta}{\pi\rho}+\frac{S^4}{n^2}\right)\Id_\ell.
\end{equation}
The relative correction inside the logarithmic entropy bound satisfies
\[
 \log\left(1+\frac{\pi\rho S^4}{n^2\eta}\right)
 =O\bigl(n^{-2+\alpha+o(1)}\bigr).
\]
The maximum-entropy bound by covariance and the fact that
conditioning lowers entropy thus give
\[
 \mathcal H(R_{\rm res}(F_{\rm star})\mid
                    \mathcal B(F_{\rm star}))
 \le \frac{\ell}{2}\log\frac{2\pi{\rm e}\eta}{\pi\rho}
          +O(n^{-c}).
\]

By the typical-star estimate at \(E+i\eta\) and the tree covariance
in Lemma~\ref{lem:tree-identities}, the covariance of the noise
\(\mathcal B(Y_{\rm star})\) differs from \(\Id_D\) by a polynomially
small error outside a polynomially small fraction of centers.
Gaussian total variation after translation, followed by the
entropy-continuity argument in Lemma~\ref{ent:star-edge}, allows this
noise to be replaced by \(G\) at polynomial entropy cost.

Similarly, the edge coordinates have the same entropy, up to
polynomial error, as the observations of \(\epsilon B_V+G\)
in directions \(w,v_i\). The deterministic discrepancy in the
neighbor coordinate is \(O(n^{-1+o(1)})\), because the neighbor
average is \(p\) times the center up to that error.
The map from orthonormal coordinates on \(U_i\) to these two
edge coordinates has absolute determinant \((1-p^2)^{\ell/2}\).
It follows that the left side of \eqref{eq:ent-star-edge} is at most
the left side of \eqref{eq:ent-projected-entropy} plus
\[
 \ell\left(
  \frac12\log\frac{2\pi{\rm e}\eta}{\pi\rho}
  -\frac12\log d+\frac{d-1}{2}\log a
                   -\frac d4\log(1-p^2)\right)+O(n^{-c}).
\]
Combine \eqref{eq:ent-star-edge},
\eqref{eq:ent-global-entropy}, and the potential identity
\eqref{eq:tree-potential}. Their constants agree exactly, proving
\eqref{eq:ent-projected-entropy}.

For the remaining assertions, put
\(\Lambda=\diag(\lambda_{i(1)},\ldots,\lambda_{i(\ell)})\).
Orthogonality of the normalized eigenvectors, the eigenvector
equations, and uniform permutation of neighbors give the exact
second-moment identities
\[
 \EE ZZ^\top=\EE X_iX_i^\top=\Id_\ell,\qquad
 \EE ZX_i^\top=\Lambda/d,\qquad
 \EE X_iX_j^\top=\frac{\Lambda^2-d\Id_\ell}{d(d-1)}
 \quad(i\ne j).
\]
The mean is zero because these eigenvectors are orthogonal to
constants. These identities imply
\(\EE B_VB_V^\top=\Id_D+O(n^{-1+o(1)})\), using
\(\Lambda=E\Id_\ell+O(n^{-1+o(1)})\).
Finally an actual uniform oriented edge is invariant under reversal.
Its two coordinates differ from the \(w,v_i\) observations of \(B_V\)
by \(O(n^{-1+o(1)})\). This error remains of that order in every
fixed moment because all coordinates are \(n^{o(1)}\).
The coordinate map is invertible since \(|p|<1\), so reversal
is precisely the asserted orthogonal swap on \(U_i\).
\end{proof}

We have obtained the entropy inequality, the approximate reversal
symmetries, and the second moments needed to identify the edge law.
The exact eigenvector normalization has supplied variance one; it is
now the higher moments that remain to be controlled.

To see the role of the next lemma, write
\(\varphi_L\) for the standard Gaussian density on a Euclidean space
\(L\), and define relative entropy by
\[
 \Ent(\Law(Z)\mid\varphi_L)
 =\int_L f\log(f/\varphi_L),
\]
when \(Z\) has density \(f\). For a centered vector \(Z\) in \(D\)
with identity covariance and finite differential entropies for itself
and its projections, the Gaussian entropy formula gives
\begin{equation}\label{eq:ent-deficit-sign}
 \mathcal H(Z)-\frac12\sum_{i=1}^d\mathcal H(P_iZ)
 =-\Ent(\Law(Z)\mid\varphi_D)
   +\frac12\sum_{i=1}^d
        \Ent(\Law(P_iZ)\mid\varphi_{U_i}).
\end{equation}
The constants cancel because
\(\dim D=\frac12\sum_i\dim U_i=d\ell\).
The same identity holds up to the second-moment error for the normalized
smoothing of \(B_V\). Thus the entropy bound limits the full relative
entropy by half the sum of its projected relative entropies.
The next lemma shows that the first nonzero Hermite coefficient would
violate this bound. Its strictness is where \(d\ge3\) and the open bulk
condition enter the argument.

\subsection{A strict tensor inequality}

For an integer \(j\), use the Euclidean norm on symmetric tensors
in \(D^{\otimes j}\). Let \(\sigma_i\) denote the orthogonal swap of
\(w,v_i\) on \(U_i\), with the identity on colors.

\begin{lemma}[Tensor coercivity]\label{ent:tensor}
For every fixed \(j\ge3\), there exists \(c_j>0\) such that every
symmetric \(j\)-tensor \(H\) on \(D\) satisfying
\(\sigma_i^{\otimes j}P_i^{\otimes j}H=P_i^{\otimes j}H\) for all
\(i\) obeys
\begin{equation}\label{eq:ent-tensor}
 \|H\|^2-\frac12\sum_{i=1}^d\|P_i^{\otimes j}H\|^2
 \ge c_j\|H\|^2 .
\end{equation}
\end{lemma}

\begin{proof}
We use the weighted-frame geometry of
\citet[Lemma~11.3]{BackhauszSzegedy}, with the numbering of the
published version, and then prove the tensor strictness needed here.
In the spatial plane spanned by \(w,v_i\), define the orthonormal
bisecting directions
\[
 e_i^+=\frac{w+v_i}{\sqrt{2(1+p)}},\qquad
 e_i^-=\frac{w-v_i}{\sqrt{2(1-p)}},\qquad
 s_i=\frac{e_i^++e_i^-}{\sqrt2},\qquad
 t_i=\frac{e_i^+-e_i^-}{\sqrt2}.
\]
The swap fixes \(e_i^+\), negates \(e_i^-\), and hence interchanges
the orthonormal directions \(s_i,t_i\). Their squared central
components are
\[
 \frac{1+\sqrt{1-p^2}}2,\qquad
 \frac{1-\sqrt{1-p^2}}2.
\]
Writing \(q_0=\sqrt{1-p^2}\), the strict bulk condition is
\[
 p^2<\frac{4(d-1)}{d^2}
 \quad\Longleftrightarrow\quad q_0>\frac{d-2}{d}.
\]
For \(d\ge3\), the two squared central components therefore
strictly bracket \(1/d\). The weights
\[
 b_s=\frac{q_0-(d-2)/d}{2q_0},\qquad
 b_t=\frac{q_0+(d-2)/d}{2q_0}
\]
are positive, sum to one, and satisfy
\begin{equation}\label{eq:ent-frame}
 \sum_{i=1}^d\bigl(b_s s_is_i^\top+b_t t_it_i^\top\bigr)
 =\Id_{\R\oplus{\bf1}^{\perp}}.
\end{equation}
Indeed its central component is one by the choice of weights.
Its mixed block vanishes because \(\sum_iPe_i=0\).
The remaining block is a scalar identity on \({\bf1}^{\perp}\);
its trace determines that scalar to be one.

Tensor \eqref{eq:ent-frame} with the color identity, and use it
in the first slot of \(\|H\|^2\). In each summand project the other
\(j-1\) slots onto \(U_i\). This can only decrease the value.
After projection, swap invariance makes the squared norms associated
with \(s_i\) and \(t_i\) equal. Since their weights sum to one,
the resulting value is
\(\frac12\sum_i\|P_i^{\otimes j}H\|^2\).
This proves nonnegativity of the left side of
\eqref{eq:ent-tensor}.

Suppose equality holds. Positivity of both weights implies that
every discarded squared norm vanishes. If \(R_{i,s}\) and
\(R_{i,t}\) are the projections onto
\(s_i\otimes\R^\ell\) and \(t_i\otimes\R^\ell\), then
\[
 R_{i,s}^{(1)}H
 =R_{i,s}^{(1)}P_i^{(2)}\cdots P_i^{(j)}H,
\]
and likewise for \(t\); superscripts specify tensor slots.
Adding the two identities and using symmetry gives
\begin{equation}\label{eq:ent-slot-transfer}
 P_i^{(a)}H=P_i^{\otimes j}H
 \quad\text{for every slot }a.
\end{equation}
In particular each projection may be moved from one slot to
another when acting on \(H\).
For \(i\ne k\), the availability of three slots now gives
\begin{align*}
 P_i^{(1)}P_k^{(1)}H
 &=P_i^{(1)}P_k^{(2)}H
  =P_k^{(2)}P_i^{(3)}H\\
 &=P_i^{(3)}P_k^{(1)}H
  =P_k^{(1)}P_i^{(1)}H .
\end{align*}
Thus \([P_i,P_k]\) annihilates the first slot of \(H\).

Let \(a_i'\) be the unit vector proportional to \(Pe_i\).
The spatial part of this commutator is
\[
 \ip{a_i'}{a_k'}
       \bigl(a_i'(a_k')^\top-a_k'(a_i')^\top\bigr),
 \qquad \ip{a_i'}{a_k'}=-\frac1{d-1}.
\]
For \(d\ge3\), it is invertible on
\(\spanop(Pe_i,Pe_k)\) and vanishes on its orthogonal complement.
The simultaneous kernel of all these commutators is therefore
the central direction, tensored with the color space.
Symmetry leaves \(H\) supported in
\((w\otimes\R^\ell)^{\otimes j}\).
But its projected swap invariance would also support it in
\((v_i\otimes\R^\ell)^{\otimes j}\); the two spatial lines are
distinct, so \(H=0\).
The continuous quadratic form is consequently strictly positive
on the unit sphere of the finite-dimensional constraint space.
Compactness proves \eqref{eq:ent-tensor}.
\end{proof}

\subsection{A finite-dimensional Gaussian criterion}

We isolate the remaining inference from the graph argument. The spaces
\(D,U_i\), projections \(P_i\), and swaps \(\sigma_i\) retain the
fixed geometry defined above.

\begin{lemma}[Entropy criterion for Gaussian moments]\label{ent:gaussian-criterion}
Fix an integer \(q\ge2\). Let \(B_n\) be random vectors in \(D\)
with \(\|B_n\|\le K_n\) almost surely, where \(K_n=n^{o(1)}\) is
deterministic. Suppose that for some \(c_*>0\),
\[
 \|\EE B_n\|+\|\EE B_nB_n^{\mathsf T}-\Id_D\|=O(n^{-c_*}),
\]
and, for every \(1\le j\le q\) and every \(i\),
\[
 \left\|(\sigma_i^{\otimes j}-\Id)
              \EE(P_iB_n)^{\otimes j}\right\|=O(n^{-c_*}).
\]
Let \(G\) be an independent standard Gaussian in \(D\). Suppose also
that, uniformly for \(0\le\epsilon\le1\),
\[
 \mathcal H(\epsilon B_n+G)
 -\frac12\sum_{i=1}^d\mathcal H(P_i(\epsilon B_n+G))
 \ge-O(n^{-c_*}).
\]
Then every fixed polynomial moment of \(B_n\) of degree at most
\(q\) differs from its standard Gaussian value by \(O(n^{-c})\),
for some \(c>0\). The conclusion is uniform over families satisfying
these hypotheses with the same \(c_*\), error constants, and bound
\(K_n\).
\end{lemma}

\begin{proof}
For \(y\in D\), define the Hermite tensors \(\operatorname{He}_j(y)\)
by the generating identity
\[
 \exp\bigl(\ip{y}{t}-\|t\|^2/2\bigr)
 =\sum_{j=0}^\infty
       \frac{\ip{\operatorname{He}_j(y)}{t^{\otimes j}}}{j!},
 \qquad t\in D.
\]
Set \(H_j=\EE\operatorname{He}_j(B_n)\). We prove by induction that
\begin{equation}\label{eq:ent-hermite-small}
 \|H_j\|=O(n^{-c_j})
 \quad(1\le j\le q).
\end{equation}
This suffices because Hermite polynomials form a basis of the ordinary
polynomials. The first- and second-moment hypotheses give the cases
\(j=1,2\).

Fix \(j\ge3\), assuming the preceding orders have been proved.
Choose \(\epsilon=n^{-\tau}\), where \(\tau>0\) satisfies
\[
 2j\tau<c_*/2,\qquad
 (j-k)\tau<c_k/2\quad(1\le k<j).
\]
This choice may depend on \(j\).
Put \(u=\epsilon/\sqrt{1+\epsilon^2}\).
Dividing every entropy argument in the hypothesis by
\(\sqrt{1+\epsilon^2}\) leaves that inequality unchanged, because
the dimensions cancel. The normalized random vector is
\[
 Z_u=uB_n+\sqrt{1-u^2}\,G.
\]
Its second-moment matrix is \(\Id_D+u^2O(n^{-c_*})\).
Consequently \eqref{eq:ent-deficit-sign}, with this second-moment
error included, turns the hypothesis into
\begin{equation}\label{eq:ent-relative-deficit}
 \Ent(\Law(Z_u)\mid\varphi_D)
 -\frac12\sum_{i=1}^d
       \Ent(\Law(P_iZ_u)\mid\varphi_{U_i})
 \le O(n^{-c_*}).
\end{equation}

We claim that
\begin{equation}\label{eq:ent-relative-expansion}
 \Ent(\Law(Z_u)\mid\varphi_D)
 =u^{2j}\left(\frac{\|H_j\|^2}{2j!}+O(n^{-c'})\right),
\end{equation}
and that the analogous formula for \(P_iZ_u\) has
\(P_i^{\otimes j}H_j\) in place of \(H_j\).
Let \(f_u\) denote the density ratio of \(Z_u\) to \(\varphi_D\).
The coefficient of \(u^k\) in its Taylor expansion is
\[
 \frac1{k!}\ip{H_k}{\operatorname{He}_k(y)}.
\]
This follows by integrating against the generating exponential:
the resulting transform is
\(\EE\exp(u\ip{B_n}{t}-u^2\|t\|^2/2)\).
Pairing two generating exponentials under Gaussian measure gives
Hermite orthogonality and the identity
\[
 \left\|\frac{\ip{H_k}{\operatorname{He}_k}}{k!}\right\|_{L^2(\varphi_D)}^2
 =\frac{\|H_k\|^2}{k!}.
\]

The Taylor remainder through degree \(j\) is
\(O(u^{j+1}n^{o(1)})\) in \(L^3(\varphi_D)\).
Indeed, at a fixed source value \(b_0\), the ratio at parameter \(s\) is
\[
 (1-s^2)^{-\dim D/2}
 \exp\left(\frac{\|y\|^2}{2}
       -\frac{\|y-sb_0\|^2}{2(1-s^2)}\right).
\]
For \(0\le s\le u\) its required derivatives are bounded by
polynomials in \(\|y\|,\|b_0\|\) times
\(C\exp(Cs\|b_0\|\|y\|)\).
Since \(\|b_0\|\le n^{o(1)}\) and \(u=n^{-\tau+o(1)}\),
their \(L^3(\varphi_D)\) norms are \(n^{o(1)}\), uniformly over
the source values. Taylor's theorem and averaging prove the bound.

For a density ratio \(1+h\), the pointwise remainder after
subtracting \(h+h^2/2\) from \((1+h)\log(1+h)\) is
bounded by \(C|h|^3\) for the entire range \(h\ge-1\).
As \(\int h\,\dd\varphi_D=0\),
\[
 \Ent((1+h)\varphi_D\mid\varphi_D)
 =\frac12\int h^2\,\dd\varphi_D
       +O\left(\int|h|^3\,\dd\varphi_D\right).
\]
All lower \(H_k\), \(1\le k<j\), are polynomially small by
induction, whereas every fixed \(H_k\) is at most \(n^{o(1)}\).
The inequalities defining \(\tau\) make the lower-degree part,
together with the Taylor remainder, of size
\(u^j O(n^{-\delta})\) in \(L^3\) for some \(\delta>0\).
In particular \(\|f_u-1\|_{L^3}=u^j n^{o(1)}\), so the cubic
entropy error divided by \(u^{2j}\) is
\(O(u^j n^{o(1)})\), again polynomially small. Orthogonality
now proves \eqref{eq:ent-relative-expansion}. The same reasoning
applies to each projection.

Substitute these expansions into \eqref{eq:ent-relative-deficit}
and divide by \(u^{2j}/(2j!)\). The choice of \(\tau\) gives
\begin{equation}\label{eq:ent-hermite-deficit}
 \|H_j\|^2-\frac12\sum_{i=1}^d
              \|P_i^{\otimes j}H_j\|^2
 \le O(n^{-c''}).
\end{equation}
Hermite tensors commute with orthogonal projections and are
equivariant under orthogonal maps, so
\[
 P_i^{\otimes j}H_j
 =\EE\operatorname{He}_j(P_iB_n)
\]
in the Gaussian space \(U_i\). Since \(\operatorname{He}_j\) is a
fixed polynomial, the assumed moment symmetries give the swap
constraints of Lemma~\ref{ent:tensor} with polynomially small errors.
In fixed finite dimensions the distance to the common kernel of a
collection of linear maps is bounded by a constant times the sizes
of their images. Project \(H_j\) onto that kernel. This changes it
by a polynomially small amount, and changes the quadratic expression
in \eqref{eq:ent-hermite-deficit} by a polynomially small amount,
since \(\|H_j\|=n^{o(1)}\). Lemma~\ref{ent:tensor} therefore proves
\eqref{eq:ent-hermite-small} and closes the induction. Each step uses
only the common constants in the hypotheses, proving uniformity.
\end{proof}

\begin{proof}[Proof of Proposition~\ref{prop:edge-moments}]
Condition on a graph in the common events of
Lemmas~\ref{ent:type-count} and \ref{ent:projected-entropy} and
Proposition~\ref{prop:spectral-input}. Apply
Lemma~\ref{ent:gaussian-criterion}, with \(q\) replaced by
\(\max(q,2)\), to \(B_n=B_V\).
Complete delocalization gives its deterministic subpolynomial bound.
Lemma~\ref{ent:projected-entropy} supplies its first two moments and
all fixed-order projected moment symmetries. Averaging
\eqref{eq:ent-projected-entropy} over \(i\) gives the criterion's
entropy hypothesis. All these inputs are uniform over the selected
eigenindices. Hence every moment of \(B_V\) of degree at most \(q\)
differs polynomially little from the corresponding standard Gaussian moment
in \(D\).

The observations of that Gaussian in directions \(w,v_i\) have
covariance \(C_\ell\). The uniformly \(O(n^{-1+o(1)})\) eigenvector
residual transfers the moment estimate to the actual edge values
\((V(v),V(v'))\). For \(k\) conditionally independent edge locations,
a coordinate monomial factors across those locations. The one-edge
estimates give the asserted joint moment estimate, with independent
edge samples and, equivalently, independent \(N(0,C_k)\) vectors
across the \(\ell\) eigenindices.

The common graph events do not depend on the selected eigenindices:
Lemma~\ref{ent:type-count} applies simultaneously to every quantized
labeling, and the spectral and eigenvector residual bounds are
uniform throughout the band. The Gaussian criterion preserves this
uniformity at every fixed moment order. This proves the proposition.
\end{proof}
